\documentclass[12pt]{amsart}
%\usepackage[margin=1in]{geometry}
\pagestyle{plain}

\usepackage{amsmath,amssymb}
\usepackage{enumerate}

\usepackage[shortlabels]{enumitem}
% Set spacing for all enumerate environments
\setlist[enumerate]{itemsep=6pt, topsep=6pt, parsep=0pt}

\usepackage{graphicx}

\usepackage[left=0.75in,right=1in,bottom=0.9in]{geometry}  
% \usepackage{mathptmx}      % use Times fonts if available on your TeX system
\usepackage[T1]{fontenc}

\usepackage[parfill]{parskip}
\renewcommand{\baselinestretch}{1}

\usepackage{fancyhdr}
\pagestyle{fancy}
\fancyhf{}
\renewcommand{\headrulewidth}{0pt}

% Header with NAME and ID on first page only
\fancypagestyle{firstpage}{
    \fancyhf{}
    \fancyhead[L]{\small NAME: \underline{\hspace{3in}}}
    \fancyhead[R]{\small ID: \underline{\hspace{2in}}}
    \renewcommand{\headrulewidth}{0pt}
}

%%%%%%%%%%%%%%%%%%%%%%%%%%%%%%%%%
% SECTION DIVIDERS
%%%%%%%%%%%%%%%%%%%%%%%%%%%%%%%%%
\newcommand{\divider}{
  \vspace{-1em} % Add space above
  \noindent
  \raisebox{-0.15ex}{\textbullet}%
  \hspace{0.5em}%
  \rule[0.5ex]{0.95\textwidth}{1.0pt}%
  \hspace{0.5em}%
  \raisebox{-0.15ex}{\textbullet}%
  \vspace{0em} % Adjust space below
}

\title{ESE 2030 QUIZZAM 2 : Spring 2026}
\author{prof-g}


\begin{document}
\thispagestyle{firstpage}

\maketitle

\begin{center}
    {\bf INSTRUCTIONS}
\end{center}

\begin{itemize}
\item 
No books, notes, or calculators. Use a writing utensil and logic. 

\item 
Please follow question instructions and fill in the bubble-sheet. 

\item 
Select one answer per problem.

\item 
Each problem has one best answer. In some problems, a wrong choice may be worth partial credit. 

\item 
Be careful to fill in the correct problem number. 

\item 
These questions are written carefully: if you detect an error, read closely and do your best.

\item 
If anything is unclear, use your best judgment. 

\item 
Cheating in any form will be dealt with severely. 

\item 
If you finish early, please turn your bubble sheet over and draw (lightly) a nice picture. 

\end{itemize}

\vspace{1in}
\begin{center}
\begin{quote}
    \large
{\em In signing my name below I attest that I will neither give nor receive unauthorized assistance on this examination.}
    \normalsize
\end{quote}
\vspace{1in}

Signature: \rule{3in}{0.4pt}
\end{center}

%%%%%%%%%%%%%%%%%%%%%%%%%%%%%%%%%%%%%%%%%%%%%%%%%%%%
\newpage
%%%%%%%%%%%%%%%%%%%%%%%%%%%%%%%%%%%%%%%%%%%%%%%%%%%%



{\bf PROBLEM 1:}
% WEEK 4
% TOPIC = similarity
%
Two engineers model the same linear transformation $T:\mathbb{R}^n\to\mathbb{R}^n$ using different coordinate systems. Engineer 1 uses basis $\mathcal{A}$ and obtains matrix $A$; Engineer 2 uses basis $\mathcal{B}$ and obtains matrix $B$. Which statement is correct?

\begin{enumerate}[(A)]
\item $A$ and $B$ have the same kernel
\item $A$ and $B$ have the same image
\item $A=B$ if the bases $\mathcal{A}$ and $\mathcal{B}$ are both orthonormal
\item $A$ and $B$ are similar
\item $A$ and $B$ are invertible
\end{enumerate}

% \textbf{Correct Answer:} D
% \textbf{Explanation:} Since both matrices represent the same transformation in different bases, $B = P^{-1}AP$ where $P$ is the change-of-basis matrix from $\mathcal{A}$ to $\mathcal{B}$. This is the definition of similarity.
% \textbf{Trick answer:} A (the abstract $\ker(T)$ is basis-independent, but $\ker(A)$ and $\ker(B)$ are sets of coordinate vectors in different bases; they are generically different subsets of $\mathbb{R}^n$)
% \textbf{Trick answer:} B (same reasoning: $\operatorname{im}(T)$ is basis-independent, but $\operatorname{im}(A)$ and $\operatorname{im}(B)$ encode it in different coordinate systems)
% \textbf{Partial credit:} A or B (correct intuition that $T$ has a single kernel/image, but conflates the abstract subspace with its coordinate representation)
% \textbf{Wrong:} C (two different orthonormal bases still generically yield different matrices; orthonormality plays no special role in similarity)
% \textbf{Wrong:} E ($T$ need not be invertible; similarity preserves rank but does not force invertibility)
%
\divider
%

{\bf PROBLEM 2:}
% WEEK 5/6
% TOPIC = QR decomposition, solving systems, orthogonal matrices, triangular systems

An $n \times n$ invertible matrix $A$ has QR decomposition $A = QR$. Using this factorization, the system $A\mathbf{x} = \mathbf{b}$ reduces to:

\begin{enumerate}[(A)]
\item Solving the lower triangular system $R^T\mathbf{x} = Q\mathbf{b}$.
\item Solving the upper triangular system $R\mathbf{x} = Q\mathbf{b}$.
\item Computing $\mathbf{x} = R^TQ^T\mathbf{b}$.
\item Computing $\mathbf{x} = QR^{-1}\mathbf{b}$.
\item Solving the upper triangular system $R\mathbf{x} = Q^T\mathbf{b}$.
\end{enumerate}

% \textbf{Correct Answer:} E
% \textbf{Explanation:} From $QR\mathbf{x} = \mathbf{b}$, multiply both sides on the left by $Q^T$. Since $Q$ is orthogonal, $Q^TQ = I$, giving $R\mathbf{x} = Q^T\mathbf{b}$. The right-hand side $Q^T\mathbf{b}$ is a single matrix--vector product (cheap), and $R$ is upper triangular, so the resulting system is solved by back-substitution (Week~1). This is the practical payoff of QR: it replaces a general linear system with a triangular one, at the cost of one orthogonal multiplication.
% \textbf{Partial credit:} B (correctly identifies that the system reduces to solving $R\mathbf{x} = (\cdots)$ by back-substitution, but uses $Q$ instead of $Q^T$ on the right-hand side. This student understands the triangular structure but has not internalized $Q^{-1} = Q^T$, the defining property of orthogonal matrices.)
% \textbf{Trick answer:} A ($R^T$ is lower triangular, not upper, and uses $Q$ instead of $Q^T$; a student who transposes the wrong factor and confuses the triangular orientation picks this)
% \textbf{Trick answer:} C ($R^TQ^T = (QR)^T = A^T$, so this computes $A^T\mathbf{b}$, not $A^{-1}\mathbf{b}$; the student who confuses transpose with inverse picks this. A good distractor because it ``looks like'' it undoes the factorization.)
% \textbf{Trick answer:} D (the order is wrong: $A^{-1} = R^{-1}Q^T$, not $QR^{-1}$. This student forgets that $(QR)^{-1} = R^{-1}Q^{-1}$, i.e., that products invert in reverse order. Moreover, explicitly computing $R^{-1}$ defeats the purpose of the factorization.)
%
\divider

{\bf PROBLEM 3:}
% WEEK 5
% TOPIC = weighted inner product, orthogonality depends on inner product
%
In $\mathbb{R}^2$, consider the weighted inner product
$$\langle \mathbf{u}, \mathbf{v} \rangle_M = \mathbf{u}^T M \mathbf{v}, \qquad M = \begin{bmatrix} 3 & 1 \\ 1 & 4 \end{bmatrix}.$$
Let $\mathbf{a} = \begin{pmatrix} 1 \\ -2 \end{pmatrix}$ and $\mathbf{b} = \begin{pmatrix} 2 \\ 1 \end{pmatrix}$. Which statement is correct?

\begin{enumerate}[(A)]
\item $\mathbf{a}$ and $\mathbf{b}$ are orthogonal under both the standard dot product and $\langle \cdot, \cdot \rangle_M$.
\item $\mathbf{a}$ and $\mathbf{b}$ are orthogonal under $\langle \cdot, \cdot \rangle_M$ but not under the standard dot product.
\item $\mathbf{a}$ and $\mathbf{b}$ are orthogonal under the standard dot product but not under $\langle \cdot, \cdot \rangle_M$.
\item $\mathbf{a}$ and $\mathbf{b}$ are not orthogonal under either inner product, but $\|\mathbf{a}\|_M = \|\mathbf{a}\|$.
\item None of the above.
\end{enumerate}

% \textbf{Correct Answer:} C
% \textbf{Explanation:} Standard dot product equals $0$, so $\mathbf{a} \perp \mathbf{b}$ in the standard sense. Weighted dot product $\neq 0$, so they are \emph{not} $M$-orthogonal. The core lesson: orthogonality depends on the choice of inner product. Two vectors perpendicular in the standard geometry can fail to be perpendicular when distances are weighted differently along different axes.
% \textbf{Trick answer:} A (computes only the standard dot product and assumes all inner products agree on orthogonality)
% \textbf{Trick answer:} B (gets the conclusion backwards; a student who computes both but swaps which is zero picks this)
% \textbf{Wrong:} D (the vectors \emph{are} standard-orthogonal, contradicting the premise; also $\|\mathbf{a}\| = \sqrt{2}$ while $\|\mathbf{a}\|_M = \sqrt{10}$, so the norms differ)
% \textbf{Wrong:} E ($M$ is symmetric with positive diagonal entries, hence positive definite, so $\langle \cdot, \cdot \rangle_M$ is a valid inner product)

\divider

%%%%%%%%%%%%%%%%%%%%%%%%%%%%%%%%%%%%%%%%%%%%%%%%%%%%
\newpage
%%%%%%%%%%%%%%%%%%%%%%%%%%%%%%%%%%%%%%%%%%%%%%%%%%%%


{\bf PROBLEM 4:}
% WEEK 5
% TOPIC = Frobenius inner product, matrix inner product spaces
%
Consider the vector space $M_{m \times n}(\mathbb{R})$ of all $m \times n$ real matrices, equipped with the Frobenius inner product $\langle A, B \rangle_F = \operatorname{tr}(A^T B)$. Which of the following is TRUE?

\begin{enumerate}[(A)]
\item For square matrices, $\|A\|_F = \sqrt{\langle A, A\rangle} = \sqrt{\operatorname{tr}(A^2)}$
\item $\langle A, B \rangle_F = 0$ if and only if $A^T B = \mathbf{0}$
\item $\langle A, B \rangle_F = \sum_{i,j} a_{ij}\,b_{ij}$
\item This satisfies the definition of an inner product only for square matrices
\item None of the above
\end{enumerate}

% \textbf{Correct Answer:} C
% \textbf{Explanation:} Expanding the trace, $\operatorname{tr}(A^T B) = \sum_{j=1}^{n}\sum_{i=1}^{m} a_{ij}\,b_{ij}$, which is exactly the standard dot product applied to the $mn$ entries of $A$ and $B$ laid out as vectors. This is the key conceptual insight: the Frobenius inner product "vectorizes" a matrix and takes the ordinary dot product, so all the machinery of inner product spaces (norms, angles, orthogonality, projections) applies to matrices.
% \textbf{Trick answer:} A (confuses the trace with the Frobenius norm; $\operatorname{tr}(A)$ sums only the diagonal entries $a_{ii}$, while $\|A\|_F = \sqrt{\sum_{i,j} a_{ij}^2}$ involves \emph{all} entries. For example, $A = \begin{psmallmatrix} 0 & 3 \\ 0 & 0 \end{psmallmatrix}$ has $\operatorname{tr}(A) = 0$ but $\|A\|_F = 3$.)
% \textbf{Trick answer:} B (confuses the matrix product $AB$ with the entry-wise inner product $\langle A, B \rangle_F$; these are fundamentally different operations. The identity $I$ and the rotation $\begin{psmallmatrix} 0 & 1 \\ -1 & 0 \end{psmallmatrix}$ are Frobenius-orthogonal since $\operatorname{tr}(I^T B) = \operatorname{tr}(B) = 0$, yet $IB = B \neq \mathbf{0}$.)
% \textbf{Wrong:} D (the Frobenius inner product requires only that $A$ and $B$ have the \emph{same} dimensions; $A^T B$ is well-defined for any $A, B \in M_{m \times n}$, producing an $n \times n$ matrix whose trace is the inner product)
% \textbf{Wrong:} E (C is correct, so "None of the above" is false)
%
\divider


{\bf PROBLEM 5:}
% WEEK 5
% TOPIC = orthogonal matrices
%
Let $Q$ be an $n \times n$ orthogonal matrix. Which of the following is \emph{not} guaranteed to be true?

\begin{enumerate}[(A)]
\item $\|Q\mathbf{x}\| = \|\mathbf{x}\|$ for all $\mathbf{x} \in \mathbb{R}^n$
\item $\det(Q) = 1$
\item The rows of $Q$ form an orthonormal set
\item $Q^{-1} = Q^T$
\item $\langle Q\mathbf{x}, Q\mathbf{y} \rangle = \langle \mathbf{x}, \mathbf{y} \rangle$ for all $\mathbf{x}, \mathbf{y} \in \mathbb{R}^n$
\end{enumerate}

% \textbf{Correct Answer:} B
% \textbf{Explanation:} From $Q^TQ = I$ we get $(\det Q)^2 = 1$, so $\det(Q) = \pm 1$. Reflections, for example, have determinant $-1$. All other listed properties are guaranteed.
% \textbf{Trick answer:} C (students often memorize that the \emph{columns} are orthonormal and forget that $Q^TQ = I$ also gives $QQ^T = I$, which means rows are orthonormal too; a student unsure about rows may incorrectly flag this)
% \textbf{Partial credit:} C (shows partial understanding if student recognizes column orthonormality but is uncertain about rows)
% \textbf{Wrong:} A (norm preservation follows directly from $\|Q\mathbf{x}\|^2 = \langle Q\mathbf{x}, Q\mathbf{x}\rangle = \langle \mathbf{x}, \mathbf{x}\rangle = \|\mathbf{x}\|^2$)
% \textbf{Wrong:} D (this is the defining equivalent condition $Q^TQ = I$)
% \textbf{Wrong:} E (inner product preservation is the fundamental property of orthogonal matrices)
%
\divider

{\bf PROBLEM 6:}
% WEEK 5
% TOPIC = orthogonal complements
%
Consider $W<\mathbb{R}^7$ with $\dim W=3$. Which statement about $W^\perp$ is true?

\begin{enumerate}[(A)]
\item $W \boxplus W^\perp = \mathbb{R}^7$
\item $W \cap W^\perp = W$
\item $\dim(W^\perp) = 3$
\item $\dim\left((W^\perp)^\perp\right)=4$
\item $W$ and $W^\perp$ do not intersect
\end{enumerate}

% \textbf{Correct Answer:} A
% \textbf{Explanation:} The orthogonal decomposition theorem guarantees $\mathbb{R}^7 = W \oplus W^\perp$: every vector decomposes uniquely into a component in $W$ and a component in $W^\perp$. Since $\dim(W) = 3$, we have $\dim(W^\perp) = 7 - 3 = 4$.
% \textbf{Partial credit:} C (correctly engages with computing $\dim(W^\perp)$ but applies the wrong formula, assuming $\dim(W^\perp) = \dim(W)$ rather than $n - \dim(W)$; shows awareness that the complement has a computable dimension)
% \textbf{Wrong:} B ($W \cap W^\perp = \{\mathbf{0}\}$; a nonzero vector cannot be orthogonal to itself since $\langle \mathbf{v}, \mathbf{v} \rangle > 0$)
% \textbf{Trick answer:} C (a common error: assuming $W$ and $W^\perp$ have the same dimension; the complement formula gives $\dim(W^\perp) = 7 - 3 = 4$, not $3$)
% \textbf{Trick answer:} D ($(W^\perp)^\perp = W$, so its dimension is $3$, not $4$; students who think taking $\perp$ twice gives something new pick this)
% \textbf{Wrong:} E ($W \cap W^\perp = \{\mathbf{0}\}$, so they do intersect --- at the zero vector; saying they ``do not intersect'' is false)
\divider 

{\bf PROBLEM 7:}
% WEEK 6
% TOPIC = pseudoinverse, geometric interpretation, inner product spaces
%
Let $T: V \to W$ be a linear transformation between inner product spaces (not necessarily injective or surjective). Given $\mathbf{b} \in W$, which description of $T^\dagger \mathbf{b}$ is correct?

\begin{enumerate}[(A)]
\item Project $\mathbf{b}$ onto $(\operatorname{im} T)^\perp$, then apply $T^{-1}$ to the result.
\item Project $\mathbf{b}$ onto $\operatorname{im}(T)$; then map the projection back to the unique preimage in $(\ker T)^\perp$ via the invertible restriction $T\big|_{(\ker T)^\perp}$.
\item Apply the adjoint $T^*$ to $\mathbf{b}$; the adjoint always inverts $T$ on $\operatorname{im}(T)$.
\item Project $\mathbf{b}$ onto $\ker(T^*)$, then solve $T\mathbf{x} = \mathbf{b}$ in that subspace.
\item Find any solution to $T\mathbf{x} = \mathbf{b}$, then project $\mathbf{x}$ onto $\ker(T)$.
\end{enumerate}
%
% \textbf{Correct Answer:} B
% \textbf{Explanation:} The restriction of $T$ to $(\ker T)^\perp$ is an isomorphism onto $\operatorname{im}(T)$. The pseudoinverse exploits this in two steps: (1) project $\mathbf{b}$ onto $\operatorname{im}(T)$, discarding the component in $(\operatorname{im} T)^\perp = \ker(T^*)$; (2) apply the inverse of the restricted map $T\big|_{(\ker T)^\perp}$ to pull the projected vector back to the unique preimage in $(\ker T)^\perp$. This yields the minimum-norm least-squares solution: "minimum norm" because the output lives in $(\ker T)^\perp$, and "least squares" because the input is the closest point in $\operatorname{im}(T)$ to $\mathbf{b}$.
% \textbf{Trick answer:} A (projects onto the \emph{wrong} subspace: $(\operatorname{im} T)^\perp$ is the part of $\mathbf{b}$ that $T$ cannot reach, so there is no preimage to find there. Also, $T^{-1}$ need not exist.)
% \textbf{Trick answer:} C ($T^*$ maps $W \to V$ and sends $\operatorname{im}(T)$ to $(\ker T)^\perp$, but it does not invert $T$; it distorts magnitudes. Students who half-remember the role of $A^T$ in the formula $A^\dagger = (A^TA)^{-1}A^T$ and attribute the inversion entirely to the adjoint pick this.)
% \textbf{Trick answer:} D (projects onto $\ker(T^*) = (\operatorname{im} T)^\perp$, which is exactly the component of $\mathbf{b}$ with \emph{no} preimage under $T$. This is the opposite of what the pseudoinverse does.)
% \textbf{Trick answer:} E (projects onto the wrong subspace in the domain: the pseudoinverse produces a vector in $(\ker T)^\perp$, not $\ker(T)$. Projecting a solution onto $\ker(T)$ gives the component that $T$ sends to zero --- the opposite of the minimum-norm idea.)
%
\divider

%%%%%%%%%%%%%%%%%%%%%%%%%%%%%%%%%%%%%%%%%%%%%%%%%%%%
\newpage
%%%%%%%%%%%%%%%%%%%%%%%%%%%%%%%%%%%%%%%%%%%%%%%%%%%%



{\bf PROBLEM 8:}
% WEEK 6
% TOPIC = orthogonal projection operators
%
Let $P$ be the matrix that orthogonally projects $\mathbb{R}^n$ onto a subspace $W$ with $\dim(W) = k$, where $0 < k < n$. Which of the following is \emph{false}?

\begin{enumerate}[(A)]
\item $P^2 = P$
\item $P^T = P$
\item $\operatorname{rank}(P) = k$
\item $I - P$ is the orthogonal projection onto $W^\perp$
\item $P$ is invertible
\end{enumerate}

% \textbf{Correct Answer:} E
% \textbf{Explanation:} Since $0 < k < n$, the projection has a nontrivial kernel ($W^\perp \neq \{\mathbf{0}\}$), so $P$ is not invertible. Concretely, $P\mathbf{v} = \mathbf{0}$ for every $\mathbf{v} \in W^\perp$. All other properties are fundamental facts about orthogonal projections.
% \textbf{Partial credit:} none (all distractors are true and well-known)
% \textbf{Wrong:} A ($P^2 = P$ is the defining algebraic property of any projection: projecting twice is the same as projecting once)
% \textbf{Wrong:} B ($P^T = P$ is what makes the projection \emph{orthogonal}; it distinguishes orthogonal projections from oblique ones)
% \textbf{Wrong:} C ($\operatorname{im}(P) = W$, so $\operatorname{rank}(P) = \dim(W) = k$)
% \textbf{Wrong:} D (if $\mathbf{v} = P\mathbf{v} + (I-P)\mathbf{v}$ is the orthogonal decomposition, then $I - P$ projects onto the complementary piece $W^\perp$)
%
\divider


{\bf PROBLEM 9:}
% WEEK 5
% TOPIC = coordinates in orthonormal bases
%
Let $\mathcal{B} = \{\mathbf{q}_1, \mathbf{q}_2, \mathbf{q}_3\}$ be an orthonormal basis for $\mathbb{R}^3$. Suppose $\mathbf{v} \in \mathbb{R}^3$ satisfies
$$\langle \mathbf{v}, \mathbf{q}_1 \rangle = 4, \qquad \langle \mathbf{v}, \mathbf{q}_2 \rangle = -1, \qquad \langle \mathbf{v}, \mathbf{q}_3 \rangle = 3.$$
What is $\|\mathbf{v}\|^2$?

\begin{enumerate}[(A)]
\item $26$
\item $16$
\item $6$
\item $36$
\item Cannot be determined without knowing the vectors $\mathbf{q}_i$ explicitly
\end{enumerate}

% \textbf{Correct Answer:} A
% \textbf{Explanation:} Since $\mathcal{B}$ is orthonormal, the coordinates of $\mathbf{v}$ are the inner products: $[\mathbf{v}]_\mathcal{B} = (4, -1, 3)^T$. By Parseval's identity, $\|\mathbf{v}\|^2 = 4^2 + (-1)^2 + 3^2 = 16 + 1 + 9 = 26$.
% \textbf{Partial credit:} D (correctly identifies all three inner products as the relevant quantities for computing $\|\mathbf{v}\|^2$ via Parseval's identity, but computes $(4 + (-1) + 3)^2 = 36$ instead of $4^2 + (-1)^2 + 3^2 = 26$; the ``squares the sum vs.\ sums the squares'' error, with the underlying concept intact)
% \textbf{Trick answer:} B (computes $4^2 = 16$ and stops, or squares only the largest coordinate)
% \textbf{Trick answer:} C (adds coordinates: $4 + (-1) + 3 = 6$; confuses norm squared with the sum of coordinates)
% \textbf{Trick answer:} D (computes $(4 + (-1) + 3)^2 = 36$; squares the sum instead of summing the squares)
% \textbf{Trick answer:} E (does not realize that orthonormality makes the inner products the coordinates, and that Parseval's identity gives the norm directly; a student who thinks explicit vectors are needed has missed the key advantage of orthonormal bases)
\divider


{\bf PROBLEM 10:}
% WEEK 6
% TOPIC = fundamental theorem of linear algebra (orthogonal version)
%
Let $A$ be a $5 \times 8$ matrix with $\operatorname{rank}(A) = 3$. Consider the four fundamental subspaces of $A$: $\ker(A)$, $\operatorname{im}(A)$, $\ker(A)^\perp$, and $\operatorname{im}(A)^\perp$. Which of the following is true?

\begin{enumerate}[(A)]
\item $\ker(A)$ and $\operatorname{im}(A)$ are orthogonal complements
\item $\ker(A)$ has dimension $3$ and is a subspace of $\mathbb{R}^5$
\item $\dim(\operatorname{im}(A)^\perp) = 3$
\item $\dim(\ker(A)^\perp) = 5$
\item A vector $\mathbf{x} \in \mathbb{R}^8$ that is orthogonal to every row of $A$ must be in $\ker(A)$
\end{enumerate}

% \textbf{Correct Answer:} E
% \textbf{Explanation:} The geometric FTLA gives $\ker(A) = (\text{row space of } A)^\perp$ in $\mathbb{R}^8$. A vector orthogonal to every row of $A$ lies in $\ker(A)$. This is the orthogonal restatement of the fact that $A\mathbf{x} = \mathbf{0}$ iff each row of $A$ dots to zero with $\mathbf{x}$.
% \textbf{Trick answer:} A ($\ker(A) \subset \mathbb{R}^8$ and $\operatorname{im}(A) \subset \mathbb{R}^5$ live in different ambient spaces, so they cannot be orthogonal complements. The correct pairs are $\ker(A) \oplus \ker(A)^\perp = \mathbb{R}^8$ and $\operatorname{im}(A) \oplus \operatorname{im}(A)^\perp = \mathbb{R}^5$.)
% \textbf{Partial credit:} None
% \textbf{Wrong:} B (two errors: $\dim(\ker(A)) = 8 - 3 = 5$, not $3$; and $\ker(A) \subset \mathbb{R}^8$, not $\mathbb{R}^5$)
% \textbf{Trick answer:} C ($\operatorname{im}(A) \subset \mathbb{R}^5$ has dimension $3$, so $\dim(\operatorname{im}(A)^\perp) = 5 - 3 = 2$, not $3$. Students who think a subspace and its complement have the same dimension pick this.)
% \textbf{Trick answer:} D ($\ker(A) \subset \mathbb{R}^8$ has dimension $5$, so $\dim(\ker(A)^\perp) = 8 - 5 = 3$, not $5$. Students who confuse the dimension of the subspace with the dimension of its complement pick this.)
\divider



%%%%%%%%%%%%%%%%%%%%%%%%%%%%%%%%%%%%%%%%%%%%%%%%%%%%
\newpage
%%%%%%%%%%%%%%%%%%%%%%%%%%%%%%%%%%%%%%%%%%%%%%%%%%%%




{\bf PROBLEM 11:}
% WEEK 4
% TOPIC = change of basis (conceptual)
%
Let $\mathcal{B}$ and $\mathcal{C}$ be two bases for $\mathbb{R}^n$, and let $P$ be the matrix satisfying $[\mathbf{v}]_\mathcal{C} = P\,[\mathbf{v}]_\mathcal{B}$ for all $\mathbf{v} \in \mathbb{R}^n$. What are the columns of $P$?

\begin{enumerate}[(A)]
\item The vectors in $\mathcal{C}$, expressed in standard coordinates
\item The vectors in $\mathcal{B}$, expressed in standard coordinates
\item The vectors in $\mathcal{B}$, expressed in $\mathcal{C}$-coordinates
\item The vectors in $\mathcal{C}$, expressed in $\mathcal{B}$-coordinates
\item The standard basis vectors, expressed in $\mathcal{C}$-coordinates
\end{enumerate}

% \textbf{Correct Answer:} C
% \textbf{Explanation:} Plugging in $\mathbf{v} = \mathbf{b}_j$ (the $j$-th vector of $\mathcal{B}$), its $\mathcal{B}$-coordinates are $\mathbf{e}_j$. So $P\mathbf{e}_j = [\mathbf{b}_j]_\mathcal{C}$: the $j$-th column of $P$ is the $j$-th basis vector of $\mathcal{B}$, written in $\mathcal{C}$-coordinates.
% \textbf{Trick answer:} D (swaps the roles of $\mathcal{B}$ and $\mathcal{C}$; this would be the columns of $P^{-1}$. Students who get the direction of conversion backwards pick this.)
% \textbf{Trick answer:} B (these are the columns of the matrix that converts from $\mathcal{B}$-coordinates to \emph{standard} coordinates, not to $\mathcal{C}$-coordinates; confuses change-of-basis with the coordinate matrix)
% \textbf{Trick answer:} A (same error as B but with the wrong basis; these columns convert from $\mathcal{C}$-coordinates to standard)
% \textbf{Wrong:} E (the standard basis plays no special role in the conversion between $\mathcal{B}$ and $\mathcal{C}$)
%
\divider

{\bf PROBLEM 12:}
% WEEK 5
% TOPIC = Gram-Schmidt and QR decomposition
%
Let $A$ be an $n \times n$ invertible matrix with columns $\mathbf{a}_1, \ldots, \mathbf{a}_n$, and let $A = QR$ be its QR decomposition. Which statement best describes the relationship between Gram--Schmidt and QR?

\begin{enumerate}[(A)]
\item The columns of $Q$ are obtained by normalizing the columns of $A$; $R$ records the original norms.
\item The columns of $Q$ are the Gram--Schmidt orthonormalization of $\mathbf{a}_1, \ldots, \mathbf{a}_n$; the entry $r_{ij}$ of $R$ is $\langle \mathbf{a}_j, \mathbf{q}_i \rangle$ for $i \leq j$.
\item The columns of $R$ are the Gram--Schmidt orthonormalization of $\mathbf{a}_1, \ldots, \mathbf{a}_n$; $Q$ records the change of basis.
\item Gram--Schmidt produces $Q$; the matrix $R = Q^T A$ is diagonal with the norms $\|\mathbf{a}_i\|$ on the diagonal.
\item Gram--Schmidt and QR solve different problems: Gram--Schmidt orthogonalizes vectors while QR factors a matrix.
\end{enumerate}

% \textbf{Correct Answer:} B
% \textbf{Explanation:} Gram--Schmidt applied to $\mathbf{a}_1, \ldots, \mathbf{a}_n$ produces orthonormal vectors $\mathbf{q}_1, \ldots, \mathbf{q}_n$, which become the columns of $Q$. Since each $\mathbf{a}_j$ is a linear combination of $\mathbf{q}_1, \ldots, \mathbf{q}_j$ (only earlier vectors), we have $\mathbf{a}_j = \sum_{i=1}^{j} r_{ij}\,\mathbf{q}_i$ with $r_{ij} = \langle \mathbf{a}_j, \mathbf{q}_i \rangle$. This is exactly $A = QR$ with $R$ upper triangular.
% \textbf{Trick answer:} A (normalizing without orthogonalizing does not produce orthonormal columns; this ignores the projection/subtraction step that is the heart of Gram--Schmidt)
% \textbf{Wrong:} C (swaps the roles of $Q$ and $R$)
% \textbf{Trick answer:} D (correctly identifies $R = Q^TA$, but $R$ is upper triangular, not diagonal; it would only be diagonal if the columns of $A$ were already orthogonal)
% \textbf{Partial credit:} D (the formula $R = Q^TA$ is correct; the error is in claiming $R$ is diagonal)
% \textbf{Wrong:} E (they are the same procedure viewed differently; QR is the matrix packaging of Gram--Schmidt)
%
\divider


{\bf PROBLEM 13:}
% WEEK 4
% TOPIC = coordinates, basis dependence of representation
%
Let $\mathcal{B}$ and $\mathcal{C}$ be two different ordered bases for $\mathbb{R}^3$. Suppose vectors $\mathbf{v}$ and $\mathbf{w}$ satisfy
$$[\mathbf{v}]_{\mathcal{B}} = [\mathbf{w}]_{\mathcal{C}} = \begin{pmatrix} 2 \\ -1 \\ 3 \end{pmatrix}.$$
Which statement is correct?

\begin{enumerate}[(A)]
\item $\mathbf{v} = \mathbf{w}$, since their coordinate representations are identical
\item $\mathbf{v}$ and $\mathbf{w}$ need not be equal
\item $\|\mathbf{v}\| = \|\mathbf{w}\|$, since the norm is determined by the coordinate entries
\item $\mathbf{v} = \mathbf{w}$ if and only if both $\mathcal{B}$ and $\mathcal{C}$ are orthonormal
\item The equation $[\mathbf{v}]_{\mathcal{B}} = [\mathbf{w}]_{\mathcal{C}}$ forces $\mathcal{B} = \mathcal{C}$
\end{enumerate}

% \textbf{Correct Answer:} B
% \textbf{Explanation:} Coordinates encode a vector relative to a specific basis. Writing $\mathcal{B} = \{\mathbf{b}_1, \mathbf{b}_2, \mathbf{b}_3\}$ and $\mathcal{C} = \{\mathbf{c}_1, \mathbf{c}_2, \mathbf{c}_3\}$, we have $\mathbf{v} = 2\mathbf{b}_1 - \mathbf{b}_2 + 3\mathbf{b}_3$ and $\mathbf{w} = 2\mathbf{c}_1 - \mathbf{c}_2 + 3\mathbf{c}_3$. Since $\mathcal{B} \neq \mathcal{C}$, these are generically different vectors in $\mathbb{R}^3$. The coordinate vector is not the vector itself; it is a representation that depends on the chosen basis.
% \textbf{Trick answer:} A (the most common misconception about coordinates: conflating a vector with its coordinate representation; students who think ``same numbers = same vector'' pick this)
% \textbf{Trick answer:} C (the norm of a coordinate vector equals the norm of the original vector only when the basis is orthonormal; for a general basis, $\|\mathbf{v}\|^2$ involves cross terms from the Gram matrix of $\mathcal{B}$, so identical coordinate entries in different bases do not guarantee equal norms)
% \textbf{Trick answer:} D (orthonormality is not the relevant condition for equality; $\mathbf{v} = \mathbf{w}$ iff $2\mathbf{b}_1 - \mathbf{b}_2 + 3\mathbf{b}_3 = 2\mathbf{c}_1 - \mathbf{c}_2 + 3\mathbf{c}_3$, which depends on the specific basis vectors, not on whether they are orthonormal)
% \textbf{Wrong:} E (reverses the logic entirely; many different pairs of vectors can share coordinate representations across different bases)
%
\divider



%%%%%%%%%%%%%%%%%%%%%%%%%%%%%%%%%%%%%%%%%%%%%%%%%%%%
\newpage
%%%%%%%%%%%%%%%%%%%%%%%%%%%%%%%%%%%%%%%%%%%%%%%%%%%%



{\bf PROBLEM 14:}
% WEEK 4
% TOPIC = matrix representation, columns encode images of basis vectors
%
Let $T: V \to V$ be a linear transformation on an $n$-dimensional vector space ($n \geq 3$), and let $\mathcal{B} = \{\mathbf{b}_1, \ldots, \mathbf{b}_n\}$ be a basis for $V$. The matrix representation $[T]_\mathcal{B}$ has the property that its third column is the zero vector.

Which conclusion follows?

\begin{enumerate}[(A)]
\item $\operatorname{coker}(T)$ contains $\mathbf{b}_3$
\item $\operatorname{im}(T)$ does not contain $\mathbf{b}_3$
\item $\ker(T)$ contains $\mathbf{b}_3$
\item $\operatorname{rank}(T) = n - 1$
\item None of the above must be true
\end{enumerate}

% \textbf{Correct Answer:} C
% \textbf{Explanation:} The $j$-th column of $[T]_\mathcal{B}$ is $[T(\mathbf{b}_j)]_\mathcal{B}$. If the third column is $\mathbf{0}$, then $[T(\mathbf{b}_3)]_\mathcal{B} = \mathbf{0}$, which means $T(\mathbf{b}_3) = \mathbf{0}$. So $\mathbf{b}_3 \in \ker(T)$, and since $\ker(T) \neq \{\mathbf{0}\}$, the transformation is not injective and therefore not invertible.
% \textbf{Trick answer:} A (that is not correct...)
% \textbf{Trick answer:} B (confuses the column space of the matrix in coordinate space with the actual image of $T$. The other columns, as $\mathcal{B}$-coordinate vectors, can have nonzero third components, meaning $T(\mathbf{b}_j)$ for $j \neq 3$ may involve $\mathbf{b}_3$. For example, if $[T]_\mathcal{B}$ has first column $(0,0,1,0,\ldots)^T$, then $T(\mathbf{b}_1) = \mathbf{b}_3 \in \operatorname{im}(T)$.)
% \textbf{Partial credit:} D (correctly recognizes rank is reduced, but overstates: one zero column gives $\operatorname{rank} \leq n-1$, not necessarily $= n-1$, since the remaining $n-1$ columns could also be linearly dependent)
% \textbf{Wrong:} E (only the third column is zero; the other columns can be nonzero, so $T$ need not be the zero map)
%
\divider

{\bf PROBLEM 15:}
% WEEK 4
% TOPIC = coordinates in non-standard bases, Lagrange basis for polynomials
%
Consider $\mathcal{P}_2$, the vector space of polynomials of degree at most 2. Consider the basis $\mathcal{L} = \{L_0, L_1, L_2\}$ associated to the nodes $x = -1, 0, 1$: each $L_i \in \mathcal{P}_2$ is the unique polynomial that equals $1$ at its own node and $0$ at the other two. Explicitly,
$$L_0(-1) = 1, \quad L_0(0) = 0, \quad L_0(1) = 0,$$
$$L_1(-1) = 0, \quad L_1(0) = 1, \quad L_1(1) = 0,$$
$$L_2(-1) = 0, \quad L_2(0) = 0, \quad L_2(1) = 1.$$

What is $[p]_\mathcal{L}$ for $p(x) = x^2 - x + 2$?
\vspace{0.15in}

\begin{enumerate}[(A)]
\begin{minipage}{\textwidth}
\item $\begin{pmatrix} 2 \\ -1 \\ 1 \end{pmatrix}$ \hfill
(B) $\begin{pmatrix} 2 \\ 2 \\ 4 \end{pmatrix}$ \hfill
(C) $\begin{pmatrix} 1 \\ -1 \\ 2 \end{pmatrix}$ \hfill
(D) $\begin{pmatrix} 4 \\ 2 \\ 2 \end{pmatrix}$ \hfill\mbox{}
\end{minipage}
\item[(E)] Cannot be determined without computing $L_0$, $L_1$, and $L_2$ explicitly
\end{enumerate}

% \textbf{Correct Answer:} D
% \textbf{Explanation:} Writing $p = c_0 L_0 + c_1 L_1 + c_2 L_2$ and evaluating at each node, the defining property gives $p(\text{node}_i) = c_i$. So the coordinates are simply the evaluations of $p$ at the nodes $-1, 0, 1$:
% $p(-1) = 1 + 1 + 2 = 4$, \quad $p(0) = 0 - 0 + 2 = 2$, \quad $p(1) = 1 - 1 + 2 = 2$.
% Thus $[p]_\mathcal{L} = (4, 2, 2)^T$. The key insight is that the Lagrange basis is designed so that coordinates \emph{are} evaluations at the nodes --- no explicit computation of the basis polynomials is needed.
% \textbf{Partial credit:} B (grasps the key insight that coordinates in the Lagrange basis are evaluations at the nodes --- the elegant consequence of the defining property --- but evaluates at the wrong node set $\{0, 1, 2\}$ instead of $\{-1, 0, 1\}$; the conceptual understanding is correct and the error is in reading the problem parameters)
% \textbf{Trick answer:} A (reads off the coefficients $2, -1, 1$ from the standard-basis representation $p(x) = 2 \cdot 1 + (-1) \cdot x + 1 \cdot x^2$; conflates the monomial basis $\{1, x, x^2\}$ with $\mathcal{L}$. This is the most common error: assuming ``coordinates'' always means the coefficients in the standard form.)
% \textbf{Trick answer:} B (understands the evaluation idea but evaluates at the nodes $0, 1, 2$ instead of $-1, 0, 1$; a student who grasps the concept but misreads the node set picks this)
% \textbf{Trick answer:} C (reads off $1, -1, 2$ from $x^2, -x, +2$ in descending-degree order; a different flavor of the monomial-basis confusion, reversing the coefficient order)
% \textbf{Trick answer:} E (misses the elegant consequence of the defining property; believes one must first derive the explicit formulas for $L_0, L_1, L_2$ and then solve a linear system. A student who picks this has not understood how the basis definition directly determines coordinates.)
%
\divider




{\bf PROBLEM 16:}
% WEEK 6
% TOPIC = pseudoinverse, left inverse, least squares
%
Let $A$ be an $m \times n$ matrix with $m > n$ and $\operatorname{rank}(A) = n$ (full column rank). The pseudoinverse is $A^\dagger = (A^TA)^{-1}A^T$. Which of the following statements is TRUE?

\begin{enumerate}[(A)]
\item $A^\dagger A = I_m$, so $A^\dagger$ is a right inverse of $A$.
\item $A^\dagger$ is an $n \times m$ matrix, and $A^\dagger A = I_n$.
\item $(A^TA)^{-1}$ does not exist because $A$ is not square.
\item $A^\dagger \mathbf{b}$ solves $A\mathbf{x} = \mathbf{b}$ exactly for every $\mathbf{b} \in \mathbb{R}^m$.
\item $A^\dagger = A^{-1}$ after discarding the extra rows of $A$.
\end{enumerate}

% \textbf{Correct Answer:} B
% \textbf{Explanation:} Since $A$ is $m \times n$, we have $A^T$ is $n \times m$, so $A^TA$ is $n \times n$. Full column rank guarantees $A^TA$ is invertible, so $(A^TA)^{-1}$ is $n \times n$. Thus $A^\dagger = (A^TA)^{-1}A^T$ is $n \times m$. Computing $A^\dagger A = (A^TA)^{-1}A^TA = I_n$: the pseudoinverse is a \emph{left} inverse of $A$.
% \textbf{Partial credit:} A (correctly recognizes that $A^\dagger$ acts as a one-sided inverse yielding an identity matrix, but confuses left with right: $A^\dagger A = I_n$, not $A A^\dagger = I_m$. Also note the dimension mismatch: $A^\dagger A$ is $n \times n$, so the result cannot be $I_m$. Shows understanding of the inverse relationship but not the directionality.)
% \textbf{Trick answer:} A (gets the side wrong: $A^\dagger A = I_n$, not $AA^\dagger = I_m$. The product $AA^\dagger$ is $m \times m$ but equals the projection onto $\operatorname{im}(A)$, not $I_m$. Students who confuse left and right inverses pick this. Note also the dimension mismatch: $I_m$ is $m \times m$ but $A^\dagger A$ is $n \times n$.)
% \textbf{Trick answer:} C (a common fear: since $A$ is not square, students worry that $A^TA$ might also be defective. But $A^TA$ is $n \times n$ and full column rank of $A$ guarantees it is invertible. This distractor catches students who do not trace the dimensions carefully.)
% \textbf{Trick answer:} D ($A\mathbf{x} = \mathbf{b}$ is overdetermined ($m > n$), so exact solutions exist only when $\mathbf{b} \in \operatorname{im}(A)$. For general $\mathbf{b}$, $A^\dagger \mathbf{b}$ gives the \emph{least squares} solution minimizing $\|A\mathbf{x} - \mathbf{b}\|$, not an exact solution. Students who conflate least squares with exact solving pick this.)
% \textbf{Wrong:} E (there is no valid operation of ``discarding rows'' to produce an inverse; the pseudoinverse accounts for the rectangular shape through the $A^TA$ Gram matrix, not by truncation)
%
\divider




%%%%%%%%%%%%%%%%%%%%%%%%%%%%%%%%%%%%%%%%%%%%%%%%%%%%
\newpage
%%%%%%%%%%%%%%%%%%%%%%%%%%%%%%%%%%%%%%%%%%%%%%%%%%%%




% THIS ONE IS A BIT REDUNDANT? TOO MUCH FROM WEEK 5? SAVE FOR FINAL EXAM
% {\bf PROBLEM:}
% % WEEK 5
% % TOPIC = orthogonal complements, subspace containment

% Let $V$ be a finite-dimensional inner product space with subspaces $W_1 < W_2 < V$. Which of the following must be true?

% \begin{enumerate}[(A)]
% \item $W_1^\perp < W_2^\perp$
% \item $W_2^\perp < W_1^\perp$
% \item $W_1^\perp \cap W_2^\perp = \{\mathbf{0}\}$
% \item $\dim(W_1^\perp) + \dim(W_2^\perp) = \dim(V)$
% \item $(W_1^\perp)^\perp = W_2$
% \end{enumerate}

% % \textbf{Correct Answer:} B
% % \textbf{Explanation:} Orthogonal complement \emph{reverses} containment. If $W_1 < W_2$, then any vector orthogonal to \emph{all} of $W_2$ is certainly orthogonal to the smaller set $W_1$. Formally: $\mathbf{v} \in W_2^\perp$ means $\langle \mathbf{v}, \mathbf{w} \rangle = 0$ for all $\mathbf{w} \in W_2$; since $W_1 < W_2$, this holds for all $\mathbf{w} \in W_1$, so $\mathbf{v} \in W_1^\perp$. The larger the subspace, the harder it is to be orthogonal to everything in it, so the complement shrinks.
% % \textbf{Trick answer:} A (the most natural first instinct --- containment should be preserved, not reversed. Students who do not think carefully about the logic of ``orthogonal to a bigger set is a stricter condition'' pick this. This is the main conceptual trap of the problem.)
% % \textbf{Wrong:} C ($W_2^\perp < W_1^\perp$, so their intersection is all of $W_2^\perp$, which is generically nontrivial unless $W_2 = V$)
% % \textbf{Trick answer:} D (this would require $\dim(W_1) + \dim(W_2) = \dim(V)$, which is not forced by containment alone. For example, in $\mathbb{R}^5$ with $W_1$ a line and $W_2$ a plane: $\dim(W_1^\perp) + \dim(W_2^\perp) = 4 + 3 = 7 \neq 5$. Students who half-remember the dimension formula $\dim(W) + \dim(W^\perp) = \dim(V)$ and try to combine the two subspaces pick this.)
% % \textbf{Wrong:} E ($(W_1^\perp)^\perp = W_1$, not $W_2$; the double complement returns the original subspace, not the larger one containing it)
% %
% \divider

{\bf PROBLEM 17:}
% WEEK 5
% TOPIC = adjoints, definition of adjoint operator
%
Let $V$ and $W$ be finite-dimensional inner product spaces with inner products $\langle \cdot, \cdot \rangle_V$ and $\langle \cdot, \cdot \rangle_W$, and let $T: V \to W$ be linear. Which of the following correctly defines the adjoint $T^*$?

\begin{enumerate}[(A)]
\item For each $\mathbf{w} \in W$, the vector $T^*(\mathbf{w})$ is the unique element of $V$ satisfying \\ $\langle T(\mathbf{v}), \mathbf{w} \rangle_W = \langle \mathbf{v}, T^*(\mathbf{w}) \rangle_V$ for every $\mathbf{v} \in V$.
\item $T^*(\mathbf{w})$ is the vector in $V$ that minimizes $\|T(\mathbf{v}) - \mathbf{w}\|_W$ over all $\mathbf{v} \in V$.
\item $T^*$ is the unique linear map $W \to V$ satisfying $T^* \circ T = I_V$.
\item $T^*(\mathbf{w})$ is the orthogonal projection of $\mathbf{w}$ onto $\operatorname{im}(T)$.
\item $T^* = T^{-1}$ when $T$ is invertible, and $T^* = (A^TA)^{-1}A^T$ otherwise, where $A$ is a matrix representing $T$.
\end{enumerate}

% \textbf{Correct Answer:} A
% \textbf{Explanation:} The adjoint is defined by ``moving $T$ across the inner product'': $\langle T(\mathbf{v}), \mathbf{w} \rangle_W = \langle \mathbf{v}, T^*(\mathbf{w}) \rangle_V$ for all $\mathbf{v} \in V$, $\mathbf{w} \in W$. The Riesz representation theorem guarantees that for each fixed $\mathbf{w}$, the functional $\mathbf{v} \mapsto \langle T(\mathbf{v}), \mathbf{w} \rangle_W$ is represented by a unique vector in $V$; that vector is $T^*(\mathbf{w})$.
% \textbf{Trick answer:} B (describes the least-squares/pseudoinverse solution $T^\dagger(\mathbf{w})$, not the adjoint; students fresh from Week 6 who conflate $T^*$ with $T^\dagger$ pick this)
% \textbf{Trick answer:} C ($T^*T = I_V$ would make $T$ an isometry; this confuses the adjoint with a left inverse and fails whenever $T$ has a nontrivial kernel)
% \textbf{Trick answer:} D (the orthogonal projection onto $\operatorname{im}(T)$ is the map $TT^\dagger$, not $T^*$; and it maps $W \to W$, not $W \to V$, so the types do not even match)
% \textbf{Trick answer:} E (conflates the adjoint with the pseudoinverse; $T^*$ always exists and equals the transpose $A^T$ in standard coordinates, while $(A^TA)^{-1}A^T$ is the pseudoinverse, a different object entirely)

\divider



{\bf PROBLEM 18:}
% WEEK 5
% TOPIC = inner product axioms
%
Which of the following defines a valid inner product on $\mathbb{R}^2$?

\begin{enumerate}[(A)]
\item $\langle \mathbf{u}, \mathbf{v} \rangle = u_1 v_1 - u_2 v_2$
\item $\langle \mathbf{u}, \mathbf{v} \rangle = u_1 v_2 + u_2 v_1$
\item $\langle \mathbf{u}, \mathbf{v} \rangle = 2u_1 v_1 + u_1 v_2 + u_2 v_1 + 3u_2 v_2$
\item $\langle \mathbf{u}, \mathbf{v} \rangle = (u_1 + u_2)(v_1 + v_2)$
\item $\langle \mathbf{u}, \mathbf{v} \rangle = u_1 v_1 + u_2 v_2 + 1$

\end{enumerate}

% \textbf{Correct Answer:} C
% \textbf{Explanation:} This corresponds to $\langle \mathbf{u}, \mathbf{v} \rangle = \mathbf{u}^T M \mathbf{v}$ with $M = \begin{bmatrix} 2 & 1 \\ 1 & 3 \end{bmatrix}$. This matrix is symmetric (so symmetry holds), and its eigenvalues are positive ($\det(M)=5>0$, $\operatorname{tr}(M)=5>0$), so $M$ is positive definite. All three axioms are satisfied.
% \textbf{Trick answer:} A (symmetric and bilinear, but fails positive definiteness: $\langle (0,1),(0,1)\rangle = -1 < 0$; students who don't check PD pick this)
% \textbf{Trick answer:} B (bilinear and symmetric, but fails positive definiteness: $\langle (1,0),(1,0)\rangle = 0$ for a nonzero vector)
% \textbf{Trick answer:} D (symmetric and bilinear, but fails positive definiteness: $\langle (1,-1),(1,-1)\rangle = 0$ for a nonzero vector; this is a rank-1 form)
% \textbf{Wrong:} E (fails $\langle \mathbf{0}, \mathbf{0} \rangle = 0$; gives $1$ instead; also fails linearity)
%
\divider



{\bf PROBLEM 19:}
% WEEK 4/5
% TOPIC = bases, linearly independent sets, Gram-Schmidt (conceptual)
%
Let $V$ be an $n$-dimensional inner product space, and let $S = \{\mathbf{v}_1, \ldots, \mathbf{v}_k\}$ be a linearly independent set in $V$ with $1 < k < n$. The vectors in $S$ are \emph{not} pairwise orthogonal. Which of the following is guaranteed?

\begin{enumerate}[(A)]
\item There exists an orthonormal set of $k$ vectors whose span equals $\operatorname{span}(S)$.
\item There exist positive scalars $\alpha_i$ such that $\{\alpha_1\mathbf{v}_1, \ldots, \alpha_k\mathbf{v}_k\}$ is orthonormal.
\item There exists an orthonormal basis for $V$ that contains every vector in $S$ (unmodified) as an element.
\item The set $S$ can be orthonormalized only after first extending it to a basis for all of $V$.
\item None of the above is guaranteed.
\end{enumerate}

% \textbf{Correct Answer:} A
% \textbf{Explanation:} Gram--Schmidt applied to $\mathbf{v}_1, \ldots, \mathbf{v}_k$ produces orthonormal vectors $\mathbf{q}_1, \ldots, \mathbf{q}_k$ satisfying $\operatorname{span}\{\mathbf{q}_1, \ldots, \mathbf{q}_j\} = \operatorname{span}\{\mathbf{v}_1, \ldots, \mathbf{v}_j\}$ for each $j$. In particular the overall spans agree. This works for any linearly independent set --- no extension to a full basis is needed, and the vectors need not already be orthogonal.
% \textbf{Trick answer:} B (rescaling each $\mathbf{v}_i$ by a positive scalar preserves its direction, so $\alpha_i\mathbf{v}_i \parallel \mathbf{v}_i$. Normalization makes each vector unit length but cannot change angles between vectors. Since the $\mathbf{v}_i$ are explicitly not pairwise orthogonal, no rescaling can make them orthonormal. This catches the ``Gram--Schmidt is just normalization'' misconception.)
% \textbf{Partial credit:} None.
% \textbf{Trick answer:} C (conflates the basis extension theorem with orthonormal extension. Since the $\mathbf{v}_i$ are not pairwise orthogonal, at least two satisfy $\langle \mathbf{v}_i, \mathbf{v}_j \rangle \neq 0$, so they cannot both belong to any orthonormal set. The student remembers ``independent vectors extend to a basis'' but forgets that orthonormal extension requires modifying the originals.)
% \textbf{Trick answer:} D (a student who thinks of Gram--Schmidt as purely a ``basis $\to$ orthonormal basis'' machine believes you must first have $n$ vectors. In fact, Gram--Schmidt works on any linearly independent set of size $k \leq n$; it orthonormalizes a subspace, not just the full space.)
% \textbf{Wrong:} E (A is guaranteed, so ``None of the above'' is false)
\divider


%%%%%%%%%%%%%%%%%%%%%%%%%%%%%%%%%%%%%%%%%%%%%%%%%%%%
\newpage
%%%%%%%%%%%%%%%%%%%%%%%%%%%%%%%%%%%%%%%%%%%%%%%%%%%%



{\bf PROBLEM 20:}
% WEEK 5
% TOPIC = adjoints, transpose, domain/codomain, injectivity/surjectivity
%
Let $T: \mathbb{R}^3 \to \mathbb{R}^2$ be defined by $T(x_1, x_2, x_3) = (x_1 + 2x_3,\; x_2 - x_3)$, with the standard dot product on both spaces. Which statement about the adjoint $T^*$ is true?

\begin{enumerate}[(A)]
\item $T^*: \mathbb{R}^2 \to \mathbb{R}^3$ is surjective.
\item $T^*: \mathbb{R}^3 \to \mathbb{R}^2$ and has the same kernel as $T$.
\item $\operatorname{im}(T^*) = \ker(T)$.
\item $T^*: \mathbb{R}^2 \to \mathbb{R}^3$ is injective.
\item $T^*$ does not exist because $T$ is not invertible.
\end{enumerate}

% \textbf{Correct Answer:} D
% \textbf{Explanation:} The matrix of $T$ is $A = \begin{bmatrix} 1 & 0 & 2 \\ 0 & 1 & -1 \end{bmatrix}$, so $T^*$ is represented by $A^T = \begin{bmatrix} 1 & 0 \\ 0 & 1 \\ 2 & -1 \end{bmatrix}$, mapping $\mathbb{R}^2 \to \mathbb{R}^3$. This $3 \times 2$ matrix has rank $2$ (both columns are independent), so $T^*$ is injective.
% \textbf{Partial credit:} C (recalls the FTLA relationship between $\operatorname{im}(T^*)$ and $\ker(T)$ but writes $\operatorname{im}(T^*) = \ker(T)$ instead of $\operatorname{im}(T^*) = \ker(T)^\perp$; the ``off by a complement'' error, showing genuine FTLA awareness with an incomplete grasp of the orthogonal structure)
% \textbf{Trick answer:} A ($T^*$ maps into $\mathbb{R}^3$ with a $2$-dimensional image, so it is not surjective; students who confuse injectivity with surjectivity or who think rank $2$ means surjective into $\mathbb{R}^3$ pick this)
% \textbf{Trick answer:} B (gets the direction of $T^*$ backwards: $T^*$ maps $\mathbb{R}^2 \to \mathbb{R}^3$, not $\mathbb{R}^3 \to \mathbb{R}^2$. Also conflates the kernels: $\ker(T)$ is $1$-dimensional in $\mathbb{R}^3$ while $\ker(T^*) = \{\mathbf{0}\}$.)
% \textbf{Trick answer:} C (tempting for students who half-remember the FTLA: the correct relationship is $\operatorname{im}(T^*) = \ker(T)^\perp$, not $\ker(T)$. Here $\ker(T)$ is $1$-dimensional while $\operatorname{im}(T^*)$ is $2$-dimensional --- they are orthogonal complements in $\mathbb{R}^3$, not equal.)
% \textbf{Wrong:} E (adjoints always exist for linear maps between inner product spaces; invertibility is irrelevant)

\divider



% I THINK THIS ONE IS TOO HARD -- COMMENTED OUT
% {\bf PROBLEM:}
% % WEEK 6
% % TOPIC = normal equations, geometric meaning of least squares

% Let $A$ be an $m \times n$ matrix with $m > n$ and full column rank. The least squares solution to $A\mathbf{x} = \mathbf{b}$ satisfies the normal equations $A^TA\hat{\mathbf{x}} = A^T\mathbf{b}$. Which statement correctly explains \emph{why} the normal equations characterize the least squares solution?

% \begin{enumerate}[(A)]
% \item $A^TA$ is the simplest invertible matrix built from $A$, so we multiply both sides of $A\mathbf{x} = \mathbf{b}$ by $A^T$ to make the system solvable.
% \item The normal equations minimize $\|A\mathbf{x}\|$ subject to $\mathbf{x}$ being as close as possible to $\mathbf{b}$.
% \item $A^T$ projects $\mathbf{b}$ onto $\operatorname{im}(A)$, so $A^T\mathbf{b}$ is the closest point in $\operatorname{im}(A)$ to $\mathbf{b}$.
% \item The normal equations are equivalent to the original system $A\mathbf{x} = \mathbf{b}$, rewritten in a form that always has a solution.
% \item None of the above.
% \end{enumerate}

% % \textbf{Correct Answer:} E
% % \textbf{Explanation:} The correct geometric explanation is: the normal equations encode the orthogonality condition $A^T(\mathbf{b} - A\hat{\mathbf{x}}) = \mathbf{0}$, which says the residual $\mathbf{b} - A\hat{\mathbf{x}}$ is orthogonal to every column of $A$, hence orthogonal to $\operatorname{im}(A)$. This is precisely the condition for $A\hat{\mathbf{x}}$ to be the closest point in $\operatorname{im}(A)$ to $\mathbf{b}$. None of (A)--(D) captures this.
% % \textbf{Trick answer:} A (multiplying by $A^T$ is not a computational trick for convenience; it encodes a geometric requirement. Students who view the normal equations as an algebraic manipulation rather than a geometric condition pick this.)
% % \textbf{Trick answer:} C (the most dangerous distractor: $A^T$ maps $\mathbb{R}^m \to \mathbb{R}^n$, so it is \emph{not} a projection onto $\operatorname{im}(A) \subset \mathbb{R}^m$. The projection is $A(A^TA)^{-1}A^T$, not $A^T$. Students who conflate the transpose with the projection pick this.)
% % \textbf{Wrong:} B (least squares minimizes $\|A\mathbf{x} - \mathbf{b}\|$, not $\|A\mathbf{x}\|$; this confuses the residual with the image)
% % \textbf{Wrong:} D (the normal equations are \emph{not} equivalent to $A\mathbf{x} = \mathbf{b}$; $A\hat{\mathbf{x}} \neq \mathbf{b}$ in general. The normal equations characterize the \emph{best approximation}, not an exact solution.)
% % \textbf{Partial credit:} C (correctly identifies that the geometry involves a closest point in $\operatorname{im}(A)$, but misidentifies $A^T$ as the operator that achieves this)

% \divider


{\bf PROBLEM 21:}
% WEEK 6
% TOPIC = orthogonal projection, error, orthogonal decomposition
%
Let $W$ be a subspace of an inner product space $V$, and let $\mathbf{v} \in V$. Define $\hat{\mathbf{v}} = \Pi_W(\mathbf{v})$ (the orthogonal projection onto $W$) and $\mathbf{r} = \mathbf{v} - \hat{\mathbf{v}}$ (the ``error''). Which of the following is guaranteed to be true?

\begin{enumerate}[(A)]
\item $\|\hat{\mathbf{v}}\| = \|\mathbf{v}\|$
\item $\langle \hat{\mathbf{v}},\, \mathbf{r} \rangle = 1$
\item $\mathbf{r} \in W$
\item $\|\mathbf{r}\| > 0$
\item None of the above
\end{enumerate}

% \textbf{Correct Answer:} E
% \textbf{Explanation:} The orthogonal decomposition gives $\mathbf{v} = \hat{\mathbf{v}} + \mathbf{r}$ with $\hat{\mathbf{v}} \in W$ and $\mathbf{r} \in W^\perp$, satisfying $\langle \hat{\mathbf{v}}, \mathbf{r} \rangle = 0$. By the Pythagorean theorem, $\|\mathbf{v}\|^2 = \|\hat{\mathbf{v}}\|^2 + \|\mathbf{r}\|^2$, so $\|\hat{\mathbf{v}}\| \leq \|\mathbf{v}\|$ with equality iff $\mathbf{v} \in W$. None of (A)--(D) is universally true.
% \textbf{Trick answer:} A (confuses projection with an isometry; projection preserves norm only when $\mathbf{v} \in W$ already, in which case $\mathbf{r} = \mathbf{0}$. For any $\mathbf{v} \notin W$, Pythagoras gives strict inequality.)
% \textbf{Trick answer:} B (remembers that the projection and residual have a special inner-product relationship, but the value is $0$, not $1$; orthogonality is the defining geometric property of the decomposition)
% \textbf{Partial credit:} C (correctly recognizes that $\mathbf{r}$ belongs to a specific subspace determined by $W$, but picks $W$ instead of $W^\perp$ --- the most common swap error in orthogonal decomposition)
% \textbf{Trick answer:} D (true for ``generic'' $\mathbf{v}$, and true in the overdetermined least-squares setting students are most familiar with, but fails when $\mathbf{v} \in W$: then $\hat{\mathbf{v}} = \mathbf{v}$ and $\mathbf{r} = \mathbf{0}$. Students who forget the edge case pick this.)
\divider


{\bf PROBLEM 22:}
% WEEK 4
% TOPIC = coordinate isomorphism, linear dependence detected via coordinates
%
Let $V$ be a 3-dimensional vector space with basis $\mathcal{B} = \{\mathbf{b}_1, \mathbf{b}_2, \mathbf{b}_3\}$ (not further specified). Three vectors in $V$ have coordinates
$$[\mathbf{u}]_\mathcal{B} = \begin{pmatrix} 1 \\ 2 \\ -1 \end{pmatrix}, \qquad [\mathbf{v}]_\mathcal{B} = \begin{pmatrix} 3 \\ 0 \\ 1 \end{pmatrix}, \qquad [\mathbf{w}]_\mathcal{B} = \begin{pmatrix} 4 \\ 2 \\ 0 \end{pmatrix}.$$

Which statement about $\{\mathbf{u}, \mathbf{v}, \mathbf{w}\}$ is correct?

\begin{enumerate}[(A)]
\item They are linearly independent because $V$ is 3-dimensional and we have exactly 3 vectors.
\item Whether they are linearly dependent or independent cannot be determined without knowing the basis $\mathcal{B}$ explicitly.
\item They are linearly independent because $\mathcal{B}$ is a basis.
\item They are linearly dependent because $[\mathbf{w}]_\mathcal{B} = [\mathbf{u}]_\mathcal{B} + [\mathbf{v}]_\mathcal{B}$.
\item None of the above.
\end{enumerate}

% \textbf{Correct Answer:} D
% \textbf{Explanation:} The coordinate map $\mathbf{v} \mapsto [\mathbf{v}]_\mathcal{B}$ is a linear isomorphism from $V$ to $\mathbb{R}^3$. Because it is linear, any linear relation among vectors is faithfully reflected in their coordinate vectors: if $[\mathbf{w}]_\mathcal{B} = [\mathbf{u}]_\mathcal{B} + [\mathbf{v}]_\mathcal{B}$, then $\mathbf{w} = \mathbf{u} + \mathbf{v}$, making the set $\{\mathbf{u}, \mathbf{v}, \mathbf{w}\}$ linearly dependent. One can verify the arithmetic: $(1,2,-1)^T + (3,0,1)^T = (4,2,0)^T$. No knowledge of $\mathcal{B}$ is needed beyond the fact that it is a basis.
% \textbf{Trick answer:} A (having $n$ vectors in an $n$-dimensional space is necessary for a basis, but not sufficient for independence; three coplanar vectors in $\mathbb{R}^3$ are dependent. This conflates a necessary condition with a sufficient one.)
% \textbf{Partial credit:} B (the student is cautious about working in an abstract space and senses that the answer depends on the relationship between $V$ and $\mathbb{R}^3$, but has not internalized the key theorem: the coordinate map is an isomorphism, so linear independence questions can always be settled from coordinates alone. This is a ``right instinct, missing theorem'' answer.)
% \textbf{Trick answer:} C (checks only for pairwise proportionality, which rules out one specific type of dependence but misses the general case. The three coordinate vectors are indeed pairwise non-proportional, yet they satisfy $[\mathbf{u}]_\mathcal{B} + [\mathbf{v}]_\mathcal{B} - [\mathbf{w}]_\mathcal{B} = \mathbf{0}$. Students who reduce ``linear independence'' to ``no vector is a multiple of another'' will pick this.)
% \textbf{Wrong:} E (D is correct, so ``None of the above'' is false)



\end{document}